% sections/related_work.tex — Related Work
\section{Related Work}
\label{sec:related_work}

We position \system{} at the intersection of four research streams:
synthetic tabular data generation, accounting-domain data synthesis,
anomaly injection for fraud detection, and privacy-preserving data sharing.

% ─────────────────────────────────────────────────────────────────────
\subsection{General-Purpose Tabular Data Generation}
\label{sec:rw_tabular}

Deep generative models have become the dominant paradigm for synthetic
tabular data.  CTGAN~\citep{xu2019modeling} introduced mode-specific
normalization and conditional training for mixed-type columns.
TVAE~\citep{xu2019modeling} applies variational autoencoders with the same
preprocessing.  Subsequent work has refined these approaches:
CopulaGAN~\citep{patki2016synthetic} combines copula transformations with
GAN training, REaLTabFormer~\citep{solatorio2023realtabformer} leverages
autoregressive transformers for relational tables, and
TabDDPM~\citep{kotelnikov2023tabddpm} adapts diffusion models to tabular
settings.

While these methods excel at preserving marginal distributions and pairwise
correlations, they treat tables as independent collections of rows and do not
enforce cross-table referential integrity, temporal ordering constraints, or
domain-specific algebraic invariants (e.g., balanced entries).  \system{}
addresses this gap by embedding accounting axioms directly into the
generation process.

% ─────────────────────────────────────────────────────────────────────
\subsection{Accounting-Domain Data Synthesis}
\label{sec:rw_accounting}

Domain-specific synthetic data for accounting has received comparatively
little attention.  The GL.ai dataset~\citep{schreyer2022federated} provides
federated synthetic journal entries but lacks document-flow linkage and
subledger coherence.  The SAFD (Simulated Accounting Fraud
Dataset)~\citep{sharma2023synthetic} generates entries with injected
anomalies, yet uses uniform random amounts that violate Benford's law.
BankSim~\citep{lopez2018banksim} focuses narrowly on payment transactions
for fraud research.

In contrast, \system{} generates a full enterprise data stack: chart of
accounts, master data, multi-process document flows with three-way matching,
subledgers, trial balances, financial statements, intercompany eliminations,
period-close entries, and OCEL~2.0 event logs---all mutually consistent.

% ─────────────────────────────────────────────────────────────────────
\subsection{Anomaly Injection and Fraud Simulation}
\label{sec:rw_anomaly}

Fraud detection benchmarks typically rely on post-hoc label injection
or manual scenario creation.  The IEEE-CIS fraud dataset~\citep{ieee_cis2019}
provides real transaction labels but lacks accounting context.
\citet{pourhabibi2020fraud} survey fraud detection methods and note the
scarcity of labeled multi-dimensional accounting fraud data.
\citet{dbouk2023anomaly} inject statistical anomalies into time series but
do not model multi-stage fraud schemes that evolve over time.

\system{} introduces a \emph{scheme-based} injection model where fraud
evolves through state-machine stages (e.g., vendor kickback schemes that
escalate over months), producing temporally correlated, multi-transaction
anomaly patterns with ground-truth difficulty and severity labels.

% ─────────────────────────────────────────────────────────────────────
\subsection{Privacy-Preserving Data Sharing}
\label{sec:rw_privacy}

Differential privacy~\citep{dwork2006calibrating} provides formal guarantees
for statistical queries.  DPGAN~\citep{xie2018differentially} and
PATE-GAN~\citep{jordon2019pate} train generative models under DP constraints
but are computationally expensive and may degrade statistical fidelity.
The synthetic data vault (SDV)~\citep{patki2016synthetic} framework supports
metadata-driven generation but does not provide formal privacy guarantees on
the metadata itself.

\system{}'s fingerprint module takes a different approach: it extracts
distributional summaries (not raw data) under $\varepsilon$-DP with explicit
$\varepsilon$-budget tracking and $k$-anonymity suppression, producing a
compact \texttt{.dsf} archive from which a fully configured generator can be
instantiated.  This decouples the privacy boundary (extraction) from the
generation process (synthesis), enabling auditability and composability.

% ─────────────────────────────────────────────────────────────────────
\subsection{Process Mining and Event Logs}
\label{sec:rw_process_mining}

Object-centric event logs (OCEL~2.0)~\citep{ghahfarokhi2021ocel} generalize
traditional single-case-ID event logs by associating events with multiple
typed objects.  While tools such as PM4Py~\citep{berti2019process} support
OCEL~2.0 import/export, generating realistic synthetic OCEL~2.0 logs with
consistent underlying business data has not been addressed.  \system{}
produces OCEL~2.0 event logs that are traceable to the same journal entries,
document flows, and master-data records that constitute the rest of the
synthetic dataset.

% ─────────────────────────────────────────────────────────────────────
\subsection{Accounting Network Construction}
\label{sec:rw_networks}

The representation of double-entry bookkeeping as directed graphs dates
back to~\citet{ijiri1965generalized}, who established that journal entries
have a natural graph interpretation: accounts are nodes, and debit/credit
flows form directed edges.  \citet{fellingham2018double} provides a modern
algebraic treatment of double-entry systems.
\citet{arya2010inferring} study the problem of inferring individual
transactions from aggregated financial statements, highlighting the
combinatorial complexity of reconstructing edge-level flows from
node-level balances.

Recent work has revived interest in accounting graphs for audit and
fraud detection.  \citet{guo2022picture} demonstrate that accounting
graph topology---node degree distributions, clustering coefficients,
and community structure---can discriminate between normal and fraudulent
bookkeeping patterns.  \citet{boersma2024complex} applies complex
network analysis to audit data, modeling interfirm relationships and
transaction patterns as multi-layer graphs.
\citet{liang2021pattern} develop pattern recognition and anomaly
detection methods directly on bookkeeping graph representations.

\citet{ivertowski2024hardware} present a hardware-accelerated approach
for constructing accounting networks from general ledger data, analyzing
155 real-world ISO~21378:2019--compliant datasets comprising
approximately 364 million journal entries with 2.4 billion line items.
Their analysis reveals key distributional characteristics of production
accounting data (see \Cref{sec:exp_realworld} for details), which
directly inform \system{}'s generation model.

While these approaches focus on \emph{analyzing} real accounting data as
graphs, none address the complementary problem of \emph{generating}
realistic synthetic accounting data that exhibits the same graph-theoretic
properties.  \system{} bridges this gap: the graph export module
(\Cref{sec:gen_graph}) produces transaction networks whose structural
characteristics are calibrated against the real-world distributions
reported in the literature.

% ─────────────────────────────────────────────────────────────────────
\subsection{Summary of Gaps}
\label{sec:rw_summary}

\Cref{tab:related_comparison} summarizes feature coverage across related
systems.

\begin{table}[t]
\centering
\caption{Feature comparison with related systems. Columns:
  \textbf{Bal}~=~balanced entries,
  \textbf{Ben}~=~Benford compliance,
  \textbf{Doc}~=~document-flow chains,
  \textbf{Sub}~=~subledger coherence,
  \textbf{FS}~=~financial statements,
  \textbf{Anom}~=~anomaly injection with labels,
  \textbf{Priv}~=~formal privacy guarantees,
  \textbf{OCEL}~=~process mining logs,
  \textbf{Graph}~=~graph export for GNNs.}
\label{tab:related_comparison}
\small
\begin{tabular}{l ccccccccc}
\toprule
\textbf{System} & \textbf{Bal} & \textbf{Ben} & \textbf{Doc} & \textbf{Sub}
  & \textbf{FS} & \textbf{Anom} & \textbf{Priv} & \textbf{OCEL}
  & \textbf{Graph} \\
\midrule
CTGAN / TVAE          & -- & -- & -- & -- & -- & -- & -- & -- & -- \\
SDV / CopulaGAN       & -- & -- & \checkmark$^*$ & -- & -- & -- & -- & -- & -- \\
REaLTabFormer         & -- & -- & \checkmark & -- & -- & -- & -- & -- & -- \\
GL.ai                 & \checkmark & -- & -- & -- & -- & \checkmark & \checkmark$^*$ & -- & -- \\
BankSim               & -- & -- & -- & -- & -- & \checkmark & -- & -- & -- \\
SAFD                  & \checkmark & -- & -- & -- & -- & \checkmark & -- & -- & -- \\
ACA / Guo et al.      & -- & -- & -- & -- & -- & -- & -- & -- & \checkmark$^*$ \\
\midrule
\textbf{\system{}}    & \checkmark & \checkmark & \checkmark & \checkmark
  & \checkmark & \checkmark & \checkmark & \checkmark & \checkmark \\
\bottomrule
\multicolumn{10}{l}{\footnotesize $^*$Partial support.}
\end{tabular}
\end{table}
